\scanpage{143}
Considérons la catégorie des espaces vectoriels \inserted{de dim finie} sur un corps $k$
(pour fixer les idées). Pour tout groupe fini $G$, soit
${\unsure{\mathcal A|A}}^G$ la catégorie des espaces vectoriels de dim finie où
$G$ opère. C'est \deleted{aussi} une catégorie avec produit
tensoriel \inserted{(\uline{exact} svt.)} et involution \inserted{(et op. $\lambda$)} (et $\mathcal A^{\{e\}}$ est la catégorie de départ).

\deleted{On a de plus un bifoncteur $\mathcal A^G\times\mathcal A^{G'}\longrightarrow\mathcal A^G$}

Cette catégorie dépend fonctoriellement de $G$ : si on a
un hom $u:G\longrightarrow G'$, on a un foncteur \uline{contravariant}
\[
\deleted{\unsure{\mathrm{Res}}}\,\mathcal A^u:\mathcal A^{G'}\longrightarrow\mathcal A^G \tag*{(1)}
\]
% slanted
\margin{$\mathcal A^u$ est un ``isomorphisme dans'' si $u$ est surjectif}
\[
\deleted{\unsure{(2)}}\,\mathcal A^{\unsure{1_G|\mathrm{id}_G}}=\deleted{\unsure{\mathcal A}}\,1_{\mathcal A^G}\ ;\quad \mathcal A^{uv}=\mathcal A^v\mathcal A^u \tag*{$(1')$.}
\]
\deleted{(le foncteur est contravariant)}. Cela permet de définir de façon purement
formelle le produit tensoriel externe
\[
\mathcal A^G\times\mathcal A^{G'}\longrightarrow\mathcal A^{G\times G'} \tag*{(2)}
\]
en posant, si $p_1,p_2$ sont les projections de $G\times G'$ sur ses
deux facteurs
\[
M\boxtimes M'=\deleted{\unsure{p_1^*(M)\otimes p_2^*(M')}}\,\mathcal A^{p_1}(M_1)\otimes\mathcal A^{p_2}(M'_2) \tag*{$(2')$}
\]
De \unsure{même} on définit la somme directe extérieure $\boxplus$
\deleted{(on peut donc dire}
% slanted
\margin{dans un sens bien compliqué à préciser !!}

$\boxtimes$ est un bifoncteur additif, \inserted{``commutatif'', ``\uline{associatif}''.} Sa connaissance
(indépendamment des structures multiplicatives dans les
$\mathcal A^G$) permet de reconstruire ces dernières, par
considération de $\Delta_G:G\longrightarrow G\times G$ : on a
\[
M\otimes M'=\deleted{\unsure{\Delta_G^*}}\,\mathcal A^{\Delta_G}(M\boxtimes M') \tag*{(3)}
\]
Enfin, on a aussi une notion de \uline{représentation}
\uline{induite}, qui pour un groupe \deleted{$G$ et un ss-groupe} \inserted{hom \uline{injectif}}
$u:G\longrightarrow G'$, définit un foncteur \uline{covariant} additif
% archive sheet number: 142
% printed form (paper watermark, vertical, cut off at both ends)
\unsure{YPEWRITER} \tikz[baseline=-.4ex,scale=.15]{\foreach \a in {0,45,...,315}{\draw (\a:.52) circle (.22);} \draw (0,0) circle (.21);} ARGENTINE \tikz[baseline=-.2ex,scale=.17]{\draw (0,0)--(1.6,0)--(1.8,.33)--(.5,.33)--(.32,.13)--(.92,.13);} \unsure{TYPEW}
